\section{Discussion}\label{sec:discussion}
\subsection{Engineering knowledge contribution and its measured boundary}
The interface turns engineering requirements into executable eligibility, commitment, source execution and evidence-lifetime checks. Its strongest supported contribution is an explicit, testable connection from a task event to the provenance and applicability of its numerical service. It distinguishes a usable late pose from an expired observation statement, a trajectory request from a union request, and a clean later version from a source history that has lost evidence status. Those differences matter to engineering consumers even when pose accuracy does not change.

The fair procedural control is decisive: it reproduces every reported numerical and authorization result. Thus the graph representation is not demonstrated to provide superior numerical decisions, risk bounds or exclusive rejection capability. The matched profile ablation is also redundant in this study. Typed dependency storage provides an inspectable common representation and selective invalidation, but these tests do not establish reduced engineering maintenance effort, operator error or development cost relative to a complete procedural implementation. Such lifecycle/usability outcomes would require their own experiment. The method is a contract-mediated localization/evidence service; it is not a newly demonstrated digital twin or general knowledge-graph reasoning algorithm.

This position also respects prior contract-based anytime control and finite-family risk calibration. Source replacement/replay connects those established ingredients to correlated revisions and event-specific evidence lifetime. The exposure identity explains a possible benefit of early publication; it neither guarantees that every eligible version is better nor eliminates the cost of installation. The current mean-work screen can be optimistic, so direct event calibration applies to the complete actual modeled execution rather than to a predicted deadline.

\subsection{When adaptation is worth its cost}
Relaxed/standard independent episodes reward the CART policy for recognizing some useful refinements. Those gains survive the new confirmation set and top-two influence diagnostic, while coming with more work than Fixed2. Under contention, time screening exchanges integrated loss for earlier availability. The common tight budget makes its additional predicate redundant. Ordinary routes frequently favour Fixed2, because modest difficulty offers little return on learned forecasting and extra publications. Strong map/geometry/prior degradation remains a numerical failure region; refusing an unsupported statement does not fix a pose.

These regions argue for retaining fixed policies inside the task-aware catalogue, not deleting inconvenient controls. The calibration selector may choose an older or Flat policy. Its mean calibration objective can also choose a trade-off that is not optimal on a finite confirmation sample. Simultaneous event bounds justify eligibility under the stated model; they do not prove that the selected policy has minimum true loss or the best hardware deployment cost.

\subsection{Reliability of the evidence}
The new protocol separates development/validation, excluded pilot characterization, calibration, fresh confirmation and unsupported stress. It preserves every seed and counts shared outcomes as paired. Numerical F0, pre-F1 claim amendment, F1 and binding records expose the actual revision chronology. The source-history amendment was made after some calibration execution and before F1; it was limited to attachment semantics, with an explicit projection theorem and conformance audit, rather than concealed as a new initial freeze. Two missing index records were recovered on their original seeds with exact raw-array/file equality. Their recovered CPU metadata is a new measurement. Local manifests are consistency evidence rather than external preregistration or protection from a malicious investigator.

Exact CP statements concern iid complete draws from fixed declared mixtures, conditional on the frozen service clock/predictors/artifacts and ideal simulator random streams. The finite family permits calibration-dependent selection, but the remaining receipt-conditioned subset, individual geometry subgroup and dependent route population have no derived conditional guarantee. Small margins, particularly the contended unfitted trajectory card and the 0.10 observation diagnostic, remain visible. Sensitivity grids are descriptive. Metadata matching certifies neither the physical sampling law nor unknown-distribution detection.

\subsection{Integration and deployment conditions}
The live implementation executes selection, precommit, retained computation, publication, replacement, causal replay and applicability checks. The guarded confirmation driver additionally executes source-history issuance in both implementations. Whole routes refuse transferred evidence and continue explicit provisional service; they do not test a physical feedback controller, online recalibration or a detector of hidden operating changes. Source scopes persist within the trusted registry lifetime; a restart requires restoration or attachment refusal. Producer/code trust, finite payload verification and current-dependency validation are required. Serialized receipt snapshots require consumer revalidation after dependency changes; no remote revocation notification is implemented.

The operation clock is a frozen mean-cost simulation, and the hard sender cap holds only in that normalized model. Added observer CPU is measured, with complete-route online costs separated from inclusive program replay and metadata microbenchmarks. Future physical deployment needs platform-specific service characterization and independent population calibration. The four original photographs and five supplied simulation views preserve engineering context without supplying missing physical performance measurements.

\subsection{What the task extension adds}
The new evidence concerns three kinds of engineering decisions. A consumer-only requirement relaxation can inherit an anchor bound when the complete failure event is pointwise covered. A producer edit can retain that bound only for a proven invariant requested projection. A remaining request requires direct renewal or refusal. The domain quantities, units, source lineage and dependency witnesses make these decisions inspectable. This is stronger than checking whether a stored card's label matches a request, but it is also implementable by a complete procedural system.

The strongest procedural comparator demonstrates that fact. It is given all dependency knowledge, event comparisons, records and caches and reproduces the contract's authorization and numerical outputs. Accordingly, the independent value shown here belongs to \emph{explicit executable semantics}, rather than a graph data structure. The interface centralizes the interpretation and produces a checkable witness; no experiment establishes that engineers maintain it faster, that a graph database is required, or that generic programming cannot achieve it. Change reasoning and interval refinement already exist in engineering informatics. Our claim is their concrete integration with finite-sample event evidence and the measured inspection trade-offs, not invention of those general tools.

Direct renewal is also a serious competitor. It requires more event-record evaluations but supports additional requests and can cache subsequent queries. Semantic closure is useful when applicable anchors already exist and repeated relaxations matter; direct renewal is useful when sufficient applicable raw outputs exist or the anchor bound is conservative. An engineering implementation should permit both. Restricting the direct baseline's old-card access would weaken that comparison; the pre-confirmation parity amendment removes that restriction and preserves its prior supported services.

The ROI experiment demonstrates selective transfer rather than complete task recovery. A geometric optical quantity can remain invariant while target coverage and delivery cost change. The full reporting service remains uncertified at the primary allowance. Relaxing risk or deleting quality requirements defines a different service; the paper reports such alternatives explicitly instead of treating them as improvements to the original request. Fixed2's retained numerical trade-offs likewise prevent an unsupported claim of universal system-level superiority.

External validity remains limited. Synthetic palette landmarks are diagnostic targets, not natural surface defects. The virtual producer uses prescribed trajectories, declared perturbations and a frozen camera-work clock. Native CPU observations are quantized and depend on a desktop environment; they do not characterize an embedded camera's real-time operation. Real deployments can contain unmodeled lighting, texture, exposure-control, scheduling and mechanical couplings that invalidate a dependency projection. A metadata match cannot reveal such changes. New matched calibration, physical evidence and a suitable distribution-monitoring mechanism would be required before claiming physical reliability.

The three task families share one geometric source and executor. They establish a broader executable consumer interface than localization alone, but do not establish generalization to unrelated engineering domains. Deterministic malformed-source fixtures test specified refusal logic, not adversarial robustness in an open system. Complete-route performance remains historical descriptive evidence without mission-level task certification. These limitations are substantive: the study improves the evidence interpretation and cross-task validation, while the representation's broad originality and application reach remain questions for peer review.


